\documentclass[11pt]{article}
\usepackage[a4paper,margin=27mm]{geometry}
\usepackage{amsmath,amssymb,amsthm,hyperref}
\emergencystretch=2em
\newtheorem{theorem}{Theorem}
\newtheorem{lemma}[theorem]{Lemma}
\title{A six-face die in a permutation-fair collection of five dice}
\author{Research note, 30 September 2026}
\date{}
\begin{document}
\maketitle

Let $g(n)$ denote the smallest number of faces that one specified die can
have in a permutation-fair collection of $n$ finite uniform dice; the
other dice may have arbitrarily many faces.

\begin{theorem}
There are five permutation-fair finite dice such that one has exactly
six faces and the other four all have the same number of faces.
Consequently, together with the four-face exclusion in
\path{no_four_face_die_for_five.tex},
\[
 5\le g(5)\le6.
\]
\end{theorem}

We first give an existence proof with a single corrected density.
An explicit rational certificate is described afterward.

\begin{lemma}[One corrected density]\label{corrector}
Let $K\ge m\ge2$, and let $u_1<\cdots<u_K$ be interior nodes of an
equal-weight quadrature rule of degree $m$ for the uniform probability
measure on $[0,1]$.
Every sufficiently close rational node vector $t_1<\cdots<t_K$ that
is still exact through degree $m-1$ admits a positive rational
piecewise-constant density $h$ with the following property:
independent variables consisting of $m-1$ uniform variables on $[0,1]$,
one variable with density $h$, and the uniform distribution on
$\{t_1,\ldots,t_K\}$ are permutation-fair.
The constant intervals of $h$ have rational endpoints, include the
$t_q$ among their boundaries, and can be chosen as the $m$ equal
subintervals of each gap between successive $t_q$, including the two
end gaps.
\end{lemma}
\begin{proof}
Call the atomic variable $D$ and the variable with density $h$ by $X$.
Fix the relative order of the $m-1$ uniform variables.
An ordering with $D<X$ is specified by nonnegative integers
$a+b+c=m-1$: there are $a$ uniform variables below $D$, $b$ between
$D$ and $X$, and $c$ above $X$. Its probability is $\int h(x)\phi(x)\,dx$,
where
\[
 \phi(x)=\frac1K\sum_{q:t_q<x}
 \frac{t_q^a(x-t_q)^b(1-x)^c}{a!b!c!}.
\]
For $X<D$, the corresponding kernels are
\[
 \phi(x)=\frac1K\sum_{q:x<t_q}
 \frac{x^a(t_q-x)^b(1-t_q)^c}{a!b!c!}.
\]
There are $m(m+1)$ kernels; all desired integrals are
$1/(m+1)!$. The other orders of the uniform labels have the same
probabilities by exchangeability.

We identify all linear relations between these kernels.
The triangular Bernstein polynomials in either display form a basis
of the bivariate polynomials of total degree at most $m-1$.
Thus a linear combination of kernels has the form
\[
 \Phi(x)=\frac1K\left(
   \sum_{t_q<x}P_+(t_q,x)+\sum_{t_q>x}P_-(t_q,x)\right),
 \qquad \deg P_\pm\le m-1.
\]
If $\Phi$ vanishes on all gaps, subtracting its polynomial expressions
on consecutive gaps gives
$P_+(t_q,x)-P_-(t_q,x)=0$ as a polynomial in $x$, for every $q$.
Since $K\ge m$, polynomial interpolation in the first variable implies
$P_+=P_-=P$ identically. Exactness of the nodes through degree $m-1$
then gives
\[
 \Phi(x)=\int_0^1P(t,x)\,dt.
\]
It follows that every relation between the kernels also annihilates
the desired right-hand sides: for a common $P$ their linear combination
is
\[
 \int_{t<x}P(t,x)\,dt\,dx+
 \int_{x<t}P(t,x)\,dt\,dx
 =\int_0^1\int_0^1P(t,x)\,dt\,dx=0.
\]
This is compatibility of the full system of ordering constraints.
Moreover its rank is constant, namely
\[
 R=m(m+1)-\left(\frac{m(m+1)}2-m\right)
   =\frac{m(m+3)}2.
\]
Here the relation space consists precisely of common polynomials $P$
whose integral in the first variable is zero; integration maps onto the
$m$-dimensional space of polynomials in $x$ of degree at most $m-1$.

Now restrict $h$ to be constant on the $m$ equal subintervals of every
gap. The resulting finite constraint matrix has the same rank $R$.
Indeed a kernel combination is polynomial of degree at most $m-1$ on
each gap. If its integral vanishes on all $m$ subintervals, it has at
least $m$ distinct zeros there, unless it vanishes identically.

At the reference nodes $u_q$, the constant density $h=1$ solves the
constraints, since these nodes are exact through degree $m$.
Choose $R$ independent rows and $R$ independent columns of the finite
constraint matrix there. The corresponding square minor stays
invertible when the nodes move slightly. Correct the constant density
only on these $R$ cells, using that square system.
Compatibility and the constant rank show that all ordering constraints
hold. The corrections tend to zero as $t\to u$, hence the resulting
density is positive. For rational $t$ the matrix and the solution are
rational. Summing all ordering probabilities gives $\int h=1$.
\end{proof}

\begin{proof}[Proof of the theorem]
Set
\[
 a_* =\left(\frac14,
 \sqrt{\frac{15/16-\sqrt{401/1280}}2},
 \sqrt{\frac{15/16+\sqrt{401/1280}}2}\right).
\]
Its coordinates are strictly increasing in $(0,1)$, and
\[
 \sum_{j=1}^3 a_{*,j}^2=1,
 \qquad \sum_{j=1}^3a_{*,j}^4=\frac35.
\]
Thus the six nodes $(1\pm a_{*,j})/2$ form an equal-weight rule
exact through degree five, in particular through degree four.

Rational points are dense near $a_*$ on the unit sphere.
For completeness, with $r=(2,3,6)/7$, the rational reflection formula
\[
 a(v)=r-2\frac{r\cdot v}{v\cdot v}v
\]
has norm one for every nonzero rational $v$, and tends to $a_*$ when
$v\to r-a_*$. Therefore there are rational $a$ arbitrarily close to
$a_*$ with $\sum a_j^2=1$. The symmetric six-node rule
$t=(1\pm a_j)/2$ is exact through degree three.
Lemma~\ref{corrector}, with $m=4$ and $K=6$, supplies three uniform
variables, one positive rational step density, and the six-face atomic
variable, all permutation-fair.

To make all variables finite, fix any equal-sided permutation-fair
four-letter block with $F$ faces per letter.
On each constant-density cell let $w_{j\ell}$ be the rational mass
of old die $j$ in that cell, and let $L$ be a common denominator of
all such masses. Replace the cell by a copy of the fair block, replacing
each occurrence of letter $j$ by $Lw_{j\ell}$ consecutive copies.
Insert a single occurrence of the distinguished letter at each of its
six cell boundaries. Each old die has $FL$ faces in total.
Block choices have the prescribed cell masses, and conditional orders
inside a cell are uniform both in the density model and in the fair
block. Hence all complete order probabilities are preserved.
Assigning distinct increasing labels to the resulting word completes
the finite construction.
\end{proof}

\paragraph{An explicit exact certificate.}
The script \path{construct_five_dice_six_faces.py} uses the reflection
direction $v=(357143,-60460,-80783)/10^7$, giving
\[
 a=\frac{(120515277394,209513296387,417091668278)}{482063494983}.
\]
Equivalently, the six distinguished nodes have common denominator
$964126989966$ and numerators
\[
\begin{gathered}
 64971826705,\quad272550198596,\quad361548217589,\\
 602578772377,\quad691576791370,\quad899155163261.
\end{gathered}
\]
The 28 constant density values are stored exactly as rational numbers
in \path{five_dice_six_face_density_certificate.json}.
They lie between $0.9995776$ and $1.0004405$.
The construction checks all 20 kernel integrals by rational arithmetic.
The independent script \path{verify_five_dice_six_faces.py} instead
uses a prefix-probability dynamic program through the intervals and
atomic jumps, checking all 120 full orders individually and exactly.
The rational cell masses have a common denominator of 1835 bits;
its size is immaterial to the six-face conclusion.

\end{document}
